% ============================================================================
% Fichier  : front/02_mode_emploi.tex
% Titre    : Mode d'emploi du cours
% Auteur   : MINKA MI NGUIDJOÏ Thierry Emmanuel
% Version  : 1.0 -- novembre 2026
% ============================================================================

\chapter*{Mode d'emploi du cours}
\addcontentsline{toc}{chapter}{Mode d'emploi du cours}
\markboth{Mode d'emploi}{Mode d'emploi}

\section*{Structure générale}

Le cours compte un prologue et six parties, soit 26~chapitres, pour
60~heures en présentiel réparties en trente séances de deux heures. Les
fiches d'auto-évaluation et les dix laboratoires de NDAP SUKULU sont des
travaux à faire hors de ces 60~heures.

\begin{tableaucours}{L{1.6cm}L{4.4cm}L{4.4cm}L{1.8cm}}{%
    \tableauHeader{Partie} &
    \tableauHeader{Thème} &
    \tableauHeader{Compétences visées} &
    \tableauHeader{Niveau}
}
Prologue & Le matin du 6 mars & Lire une crise, séparer constatation et
opinion & \niveaubase \\
I & Fondements du risque & Incertitude, biais, gouvernance, cadres, droit
& \niveaubase \\
II & Analyser les risques & Identifier, coter, EBIOS RM, FAIR,
Monte-Carlo, traiter & \niveaumoyen \\
III & Auditer le dispositif & Programme, preuves, échantillons, rapport &
\niveaumoyen \\
IV & Incidents et continuité & Incidents, analyse d'impact, plans,
reprise, exercices & \niveaumoyen \\
V & Gestion de crise & Cellule, décision, communication, exercices &
\niveaumoyen{} à \niveauexpert \\
VI & Apprendre et mesurer & Retour d'expérience, \OUI{}~/ \NON{}~/
\NEK{} & \niveauexpert \\
\end{tableaucours}

\section*{Les trente séances}

\begin{tableaucours}{L{1.4cm}L{8.4cm}L{2.4cm}}{%
    \tableauHeader{Séance} &
    \tableauHeader{Contenu} &
    \tableauHeader{Partie}
}
1 & Prologue~: le matin du 6 mars & Prologue \\
2 & Risque, incertitude et décision (Ch.~\ref{chap:risque}) & I \\
3 & Percevoir le risque (Ch.~\ref{chap:perception}) & I \\
4 & Gouvernance des risques (Ch.~\ref{chap:gouvernance}) & I \\
5 & Les cadres de management des risques (Ch.~\ref{chap:cadres}) & I \\
6 & Obligations juridiques (Ch.~\ref{chap:droit}) & I \\
7 & Identifier les risques (Ch.~\ref{chap:identifier}) & II \\
8 & Échelles et matrices (Ch.~\ref{chap:matrices}) & II \\
9--10 & ISO/IEC~27005 et EBIOS Risk Manager (Ch.~\ref{chap:ebios}) & II \\
11--12 & Analyse quantitative, FAIR (Ch.~\ref{chap:quantitatif}) & II \\
13 & Simuler~: Monte-Carlo (Ch.~\ref{chap:simulation}) & II \\
14 & Traiter, suivre, rendre compte (Ch.~\ref{chap:traiter}) & II \\
15 & L'audit selon l'ISO~19011 (Ch.~\ref{chap:audit}) & III \\
16 & Conduire l'audit (Ch.~\ref{chap:conduire}) & III \\
17 & Auditer la gestion des risques et le SMSI (Ch.~\ref{chap:auditer-risques}) & III \\
18 & Le rapport et le suivi (Ch.~\ref{chap:rapport}) & III \\
19 & Gérer les incidents (Ch.~\ref{chap:incidents}) & IV \\
20 & L'analyse d'impact sur l'activité (Ch.~\ref{chap:bia}) & IV \\
21 & Stratégies et plans de continuité (Ch.~\ref{chap:pca}) & IV \\
22 & Reprise informatique et résilience (Ch.~\ref{chap:pra}) & IV \\
23 & Exercer et tester les plans (Ch.~\ref{chap:exercer}) & IV \\
24 & Qu'est-ce qu'une crise~? (Ch.~\ref{chap:crise}) & V \\
25--26 & La cellule de crise (Ch.~\ref{chap:cellule}) & V \\
27 & Décider en crise (Ch.~\ref{chap:decider}) & V \\
28 & Communiquer en crise (Ch.~\ref{chap:communiquer}) & V \\
29 & Les exercices de crise (Ch.~\ref{chap:exercices-crise}) & V \\
30 & Retour d'expérience et mesure~; synthèse (Ch.~\ref{chap:retex}) & VI \\
\end{tableaucours}

\section*{Les encadrés}

\begin{boxdef}
\textbf{Définition}. Introduit un terme avec sa définition~; les termes
définis sont indexés et repris dans le glossaire (Annexe~\ref{ann:glossaire}).
\end{boxdef}

\begin{boxscenario}{Titre du scénario}
\textbf{Scénario opérationnel}. Un épisode d'une des opérations du
cours~: GAROUA et LIZBIZ (LiZbiZ), WOURI (BANQUE SANAGA), MOUNGO
(entrepôt de Bonabéri).
\end{boxscenario}

\begin{boxobjet}{R-xx · Intitulé}
\textbf{Fiche de risque}. Présente un risque sous une forme constante~:
scénario, cotation, traitement, propriétaire, indicateur de suivi.
\end{boxobjet}

\begin{boxprocedure}
\textbf{Procédure de crise}. Dit, pas à pas, ce qu'il faut faire.
\end{boxprocedure}

\begin{boxmetier}{Gestionnaire de risques}
\textbf{Application métier}. Traduit le chapitre en gestes concrets pour
l'un des métiers visés~: gestionnaire de risques, auditeur des systèmes
d'information, responsable de la continuité, directeur de crise, chargé de
communication de crise.
\end{boxmetier}

\begin{boxjuris}
Cite un texte de loi d'après sa version officielle, puis en tire la
portée pour la gestion des risques et des crises.
\end{boxjuris}

\begin{boxalert}
\textbf{Violation typique} et \textbf{règle}~: une erreur fréquente, et ce
qu'il faut faire à la place.
\end{boxalert}

\begin{boxexo}
\textbf{Exercice}. Dix par chapitre, gradués en trois niveaux~:
\niveaubase\ fondamental, \niveaumoyen\ intermédiaire, \niveauexpert\
expert. Les corrigés figurent dans le livret de correction.
\end{boxexo}

\section*{Conventions}

\begin{itemize}
    \item Chaque analyse indique un niveau de confiance~: faible, moyenne ou
          élevée, et reste séparée de la constatation.
    \item Chaque estimation chiffrée est donnée avec son intervalle ou ses
          hypothèses~; un chiffre seul n'est pas une estimation.
    \item Les \termecle{termes clés} apparaissent en gras~; les normes
          s'écrivent \norme{ISO~31000}, les lois camerounaises
          \loicam{n°~2010/012}.
    \item Les échelles, les seuils et les grilles proposés sont des
          conventions pédagogiques, à adapter par chaque organisation.
\end{itemize}

\section*{Matériel d'accompagnement}

\begin{itemize}
    \item \textbf{Annexes}~: glossaire, outils, modèles et procédures,
          textes juridiques et normatifs, bibliographie commentée, contacts
          (Annexes~\ref{ann:glossaire} à~\ref{ann:contacts}).
    \item \textbf{Livret de correction}~: corrigés des exercices et des
          fiches, réservé au formateur.
    \item \textbf{NDAP SUKULU}~: fiches d'auto-évaluation et laboratoires,
          présentés dans la section qui suit.
\end{itemize}
